\documentclass[lettersize,journal]{IEEEtran}
\usepackage{amsmath,amsfonts}
\usepackage{algorithmic}
\usepackage{algorithm}
\usepackage{array}
\usepackage[caption=false,font=footnotesize]{subfig}
\usepackage{textcomp}
\usepackage{stfloats}
\usepackage{url}
\usepackage{verbatim}
\usepackage{graphicx}
\usepackage{cite}
\usepackage{float}

%\usepackage{xcolor}
%\usepackage{soul}

% --- CL float stabilization (minimal) ---
\setlength{\textfloatsep}{5pt}
\setlength{\floatsep}{5pt}
\setlength{\intextsep}{5pt}

\hyphenation{op-tical net-works semi-conduc-tor IEEE-Xplore}
% updated with editorial comments 8/9/2021

\begin{document}

\title{\huge I/Q Imbalance Pre‑Compensation for OFDM Systems Using RF Feedback and BLSTM Predistortion}

\author{Jihoon~Lee and Jin~Kwan~Kim% <-this % stops a space
%\IEEEmembership{, Member,~IEEE}

\thanks{
J. Lee is with Satellite Communication Research Division, Electronics and Telecommunications Research Institute (ETRI), Daejeon 34129, Republic of Korea (e-mail: jihoonlee0104@etri.re.kr).

J. K. Kim (corresponding author) is an independent researcher, Taean, Chungcheongnam-do 32164, Republic of Korea (e-mail: tull00@hanyang.ac.kr).
}

%This work was supported by Institute for Information \& communications Technology Planning \& Evaluation (IITP) grant funded by the Korea government (KNPA) (No.2019-0-01291, LTE-based accurate positioning technique for emergency rescue).
%This work was in part supported by the research fund of Hanyang University.}
}

% The paper headers
%\markboth{Journal of \LaTeX\ Class Files,~Vol.~14, No.~8, August~2021}%
%{Shell \MakeLowercase{\textit{et al.}}: A Sample Article Using IEEEtran.cls for %IEEE Journals}

%\IEEEpubid{0000--0000/00\$00.00~\copyright~2021 IEEE}
% Remember, if you use this you must call \IEEEpubidadjcol in the second
% column for its text to clear the IEEEpubid mark.

\maketitle

\begin{abstract}
In this paper, we propose two novel predistortion schemes to mitigate I/Q imbalance in Orthogonal Frequency Division Multiplexing (OFDM) communication systems: RF signal-based and deep learning-based schemes.
With the RF signal-based scheme, gain and phase errors are estimated in real time using feedback from the RF signal.
With the deep learning-based scheme, I/Q imbalance predistortion is performed using Bi-directional Long Short-Term Memory (BLSTM). The simulation results show that the proposed schemes reduce the Bit Error Rate (BER) in OFDM systems.
\end{abstract}

\begin{IEEEkeywords}
I/Q imbalance, predistortion, deep learning, OFDM
\end{IEEEkeywords}

\section{Introduction}
\IEEEPARstart{W}{ith} the evolution from 5G New Radio (NR) to the emerging sixth generation (6G) wireless systems, there is an increasing demand for higher data rates, lower latency, and massive connectivity.
Both NR and emerging 6G systems adopt OFDM waveforms per recent 3GPP standards \cite{REF1, RP-25xxxx, RP-252912}.

OFDM offers several advantages, including high spectral efficiency and strong robustness against multi-path fading and Doppler effects \cite{Prasad}.
OFDM enables efficient FFT-based signal processing, does not require complex equalizers, and is robust to various channel impairments.
Therefore, OFDM remains widely adopted in modern wired and wireless communication systems, with extensive ongoing standardization efforts based on OFDM technologies.
However, OFDM systems are highly sensitive to impairments such as frequency offset, high Peak-to-Average Power Ratio (PAPR), DC offset, in-phase/quadrature imbalance, and phase noise \cite{REF4, REF5}.

Modern wireless systems demand highly integrated, low-cost receivers.
Zero Intermediate Frequency (zero-IF) receivers have been widely adopted \cite{REF6}.
However, while baseband processing occurs in the digital domain, up/down-conversion is performed in the analog domain, where I/Q imbalance is inevitable due to gain and phase mismatches \cite{Razavi}.
In addition, as wireless systems move toward higher frequency bands and wider bandwidths in 6G, the impact of I/Q imbalance becomes more severe, making effective compensation techniques 
increasingly critical.

Several research efforts have addressed I/Q imbalance \cite{REF7, REF8, REF9}.
However, these require additional pilot signals, reducing resource efficiency, and focus on post-compensation at the receiver rather than pre-compensation.
Since post-compensation results in larger residual errors, pre-compensation is preferable.
We propose two predistortion schemes: (1) RF signal-based without pilot signals, and (2) BLSTM-based \cite{REF10, REF11, REF12}.

The simulation results demonstrate that the proposed predistortion schemes improve the Bit Error Rate (BER) performance of the OFDM system under I/Q imbalance conditions.

\section{System model}

\subsection{OFDM System}
In an OFDM system, a high-rate data stream is divided into multiple low-rate data streams.
The time-domain OFDM signal $x(t)$ is \cite{Prasad}

\begin{equation}
\label{eq1}
x(t) = \frac{1}{N}\sum_{n=0}^{N-1} X_{n} e^{j2\pi n\Delta ft}, \quad 0 \leq t < T
\end{equation}
where $N$ denotes the number of subcarriers, $X_n$ represents the data symbol on the $n$th subcarrier, $\Delta f =1/T$ is the subcarrier spacing and $T$ is OFDM symbol duration.
In (\ref{eq1}), the time-domain OFDM signals $x(t)$ can be expressed as a sum of in-phase signal $I(t)$ and quadrature signal $Q(t)$ $(x(t) = I(t) + jQ(t))$.

\subsection{I/Q Imbalance}
In the up-conversion process, in-phase and quadrature signals are mapped onto cosine and sine waves of Local Oscillator (LO), respectively.
In ideal up-conversion, the RF signal $s(t)$ is

\begin{equation}
\label{eq2}
s(t) = I(t) \cdot cos(\omega_c t) + Q(t) \cdot sin(\omega_c t), 
\end{equation}
where $\omega_c$ is the carrier frequency of the RF signal.

%\begin{figure}[!t]
%\centering
%\includegraphics[width=0.88\linewidth]{local_oscillator}
%\caption{I/Q imbalance of the local oscillator}
%\label{lo}
%\end{figure}

%Fig. \ref{lo} shows the I/Q imbalance of the LO.
The signal with I/Q imbalance $\hat s(t)$ can be expressed as \cite{Razavi}

\begin{equation}
\label{eq3}
\begin{split}
\hat s(t) 
&= I(t) \left(1+\frac{\varepsilon}{2}\right) cos \left(\omega_c t + \frac{\varphi}{2}\right) \\
&+ Q(t) \left(1-\frac{\varepsilon}{2}\right) sin\left(\omega_c t - \frac{\varphi}{2}\right), 
\end{split}
\end{equation}
where $\varepsilon$ and $\varphi$ are the gain and phase errors, respectively.


\section{I/Q imbalance predistortion schemes}

\subsection{RF signal-based predistortion of I/Q imbalance}

\begin{figure*}[!t]
\centering
\includegraphics[width=0.88\textwidth]{Predistortor}
\caption{Transmitter structure with I/Q imbalance predistorter based on RF signal.}
\label{transmit}
\end{figure*}

Fig. \ref{transmit} shows the transmitter structure with I/Q imbalance predistorter based on the RF signal.
The RF-based scheme uses feedback from the transmitted RF signal.
A transmitter down-converts the RF signal using the LO used in the up-conversion process.
The down-conversion signal $\tilde x(t)$ can be expressed as the sum of in-phase signal $\tilde I(t)$ and quadrature $\tilde Q(t)$ $(\tilde x(t) = \tilde I(t) + j \tilde Q(t))$.
Then, $\tilde I(t)$ and $\tilde Q(t)$ can be expressed as

\begin{equation}
\label{eq4}
\begin{split}
&\tilde I(t) 
= \int_{0}^{t} \hat s(\tau) \left(1+\frac{\varepsilon}{2}\right) cos \left(\omega_c \tau + \frac{\varphi}{2}\right)+n(\tau)d\tau \\
&= I_{L}(t) \left(1+\frac{\varepsilon}{2}\right)^2 - Q_{L}(t) \left(1-\frac{\varepsilon^2}{4}\right) sin \varphi + n_{L}(t),
\end{split}
\end{equation}

\begin{equation}
\label{eq5}
\begin{split}
&\tilde Q(t) 
= \int_{0}^{t} \hat s(\tau) \left(1-\frac{\varepsilon}{2}\right) sin \left(\omega_c \tau - \frac{\varphi}{2}\right) + n(\tau)d\tau \\
&= Q_{L}(t) \left(1-\frac{\varepsilon}{2}\right)^2 - I_{L}(t) \left(1-\frac{\varepsilon^2}{4}\right) sin \varphi + n_{L}(t),
\end{split}
\end{equation}
where $I_{L}(t)$, $Q_{L}(t)$ and $n_{L}(t)$ represent the I and Q signals and internal thermal noise after low-pass filtering, respectively.

In \eqref{eq4} and \eqref{eq5}, $\varepsilon^2$ can be negligible for $0<\varepsilon < 0.1$.
Then, from (\ref{eq4}) and (\ref{eq5}), the gain error $\varepsilon_T$ and the phase error $\varphi_T$ can be estimated as

\begin{equation}
\label{eq6}
\varepsilon_T =
\frac{I(t) \tilde I(t) - Q(t) \tilde Q(t) - (I^2(t) - Q^2(t))}{I^2(t) + Q^2(t)},
\end{equation}

\begin{equation}
\label{eq7}
\varphi_T =
sin^{-1} \left( \frac{2I(t)Q(t) - (I(t) \tilde Q(t) - Q(t) \tilde I(t))}{I^2(t) + Q^2(t)} \right).
\end{equation}

\subsection{Deep learning-based predistortion of I/Q imbalance}

Under ideal conditions, the RF signal-based scheme achieves near-ideal BER performance.
However, the performance degradation of predistortion can be observed in environments with noise in the internal circuitry.
This performance degradation is considered to result from the process of converting the RF signal through the DAC and the subsequent internal feedback loop involved in the RF signal-based predistortion scheme.
In this subsection, the deep learning based predistortion of I/Q imbalance is proposed, which does not involve the internal feedback or DAC processes.

In this paper, we employ BLSTM as a deep learning-based approach for I/Q imbalance predistortion. Unlike traditional LSTM, BLSTM processes information bidirectionally, utilizing both past and future context \cite{REF10, REF11, REF12}.

At the input layer, the real and imaginary parts of the complex signal affected by I/Q imbalance are separately fed into the BLSTM model.
These signals pass through both forward and backward BLSTM layers, which learn time-dependent features.
The forward and backward hidden states, denoted by $\overrightarrow{h}_t$ and $\overleftarrow{h}_t$, are expressed as

\begin{equation}
\label{eq8}
\overrightarrow{h}_t = LSTM_{forward}(x_t,\overrightarrow{h}_{t-1}),
\end{equation}

\begin{equation}
\label{eq9}
\overleftarrow{h}_t = LSTM_{backward}(x_t,\overleftarrow{h}_{t+1}).
\end{equation}

Then, the combined hidden state $h_t$, obtained by combining the forward and backward hidden states, can be expressed as

\begin{equation}
\label{eq10}
h_t = [\overrightarrow{h}_t,\overleftarrow{h}_t].
\end{equation}

The BLSTM model is trained using pairs of signals with I/Q imbalance and their corresponding ideal signals.
The loss function $L$ is defined as the Mean Squared Error (MSE) and can be expressed as

\begin{equation}
\label{eq11}
L = \frac{1}{N} \sum_{i=1}^{N}(|I_{i,ideal}-I_{i,comp}|^2+|Q_{i,ideal}-Q_{i,comp}|^2),
\end{equation}
where $I_{i,ideal}$ and $Q_{i,ideal}$ are in-phase and quadrature part of the ideally predistorted signal, respectively.
On the other hand, $I_{i,comp}$ and $Q_{i,comp}$ are in-phase and quadrature part of the predistorted signal obtained through training, respectively.
Finally, the model is trained using the training dataset, which consists of OFDM symbols with I/Q imbalance, and the defined loss function based on the MSE. 
Through the optimization process, the model learns to effectively compensate for I/Q imbalance by minimizing the defined loss between the ideal and the compensated signals.

For I/Q imbalance predistortion, the use of BLSTM algorithm offers several advantages.
It can effectively compensate for I/Q imbalance exhibiting non-linear characteristics, which is challenging for conventional linear compensation algorithms.
Moreover, BLSTM is capable of capturing the time-dependent changes in I/Q imbalance, making it suitable for dynamically changing environments.
However, the BLSTM-based I/Q imbalance predistortion also presents certain limitations.
It requires a large amount of training set to achieve reliable performance, and the training process can be computationally intensive and time-consuming.
In addition, the implementation may demand more computational resources compared to traditional compensation techniques. 

In this paper, considering the system complexity and hardware limitations, the OFDM system is designed with an FFT size of 32.
The proposed BLSTM model consists of an input layer, two BLSTM hidden layers, and an output layer.
The first and second BLSTM hidden layers contain 70 and 90 units, respectively, chosen to balance model accuracy and computational complexity under real-time processing constraints.
The first layer captures fundamental I/Q imbalance patterns across 32 subcarriers, while the second layer models non-linear interactions between subcarriers.
This achieves near-optimal MSE performance while maintaining processing within one OFDM symbol duration.
To train the BLSTM model, the Root Mean Square Propagation (RMSProp) optimizer was used.
The batch size was set to 32, and training was conducted over 60 epochs.
The training dataset consisted of a total of 3,200,000 OFDM symbols, which were split into training (2,560,000 symbols) and test (640,000 symbols) sets at a ratio of 4:1.

\begin{figure}[!t]
\centering
\includegraphics[width=0.88\linewidth]{transmit_blstm}
\caption{Transmitter structure with I/Q imbalance predistorter by BLSTM}
\label{transmit_blstm}
\end{figure}

Fig. \ref{transmit_blstm} shows the example of OFDM transmitter structure including the BLSTM used for I/Q imbalance predistortion.
In Fig. \ref{transmit_blstm}, the BLSTM is applied after the Parallel-to-Serial (P/S) conversion. 
During the BLSTM processing, the OFDM symbols are predistorted to mitigate the I/Q imbalance and then up-converted to the RF signal.


\section{Simulation Results}

In this section, simulation results are presented. 
The OFDM signal is modulated by Quadrature Phase Shift Keying (QPSK), and the number of subcarriers is set to 32. 
For BLSTM training, the input and target data are set to signals with I/Q imbalance and the corresponding ideal signals, respectively. 
The BLSTM training is performed independently for the I and Q components. 
The I/Q imbalance is modeled by assuming a gain error uniformly distributed in the range of $0.3\sim0.6$ and a phase error uniformly distributed in the range of $5 \,^{\circ}\sim8 \,^{\circ}$.
For comparison, the RF signal-based predistortion scheme is evaluated under ideal conditions, where the effect of internal feedback loop errors is not taken into account.

The performance of the training model was evaluated using the MSE. 
On the test set, the proposed BLSTM-based predistortion model achieve
d an MSE of $1.2 \times 10^{-4}$, demonstrating stable compensation performance under the 
considered I/Q imbalance conditions.
In particular, the model demonstrated stable predistortion performance for I/Q imbalance parameters beyond those explicitly seen during training, indicating its generalization capability.


\subsection{Signal constellation under I/Q imbalance}

\begin{figure}[!t]
\centering
\subfloat[With IQ imbalance]{\label{const_iq_imbal}
    \includegraphics[width=0.8\linewidth]{const_iq_imbal}}
\vfill
\subfloat[After BLSTM predistortion]{\label{const_blstm}
    \includegraphics[width=0.8\linewidth]{const_blstm}}
\caption{Signal constellations in AWGN at $E_b/N_0 = 15$ dB.}
\label{fig:const}
\end{figure}

Fig. \ref{fig:const} shows signal constellations in AWGN at $E_b/N_0 = 15$ dB. 
In Fig.~\ref{fig:const}\subref{const_iq_imbal}, I/Q imbalance causes constellation dispersion, rotation, and multiple distinct clusters. 
This characteristic pattern arises because I/Q imbalance is modeled using discrete amplitude and phase error parameters, unlike random RF impairments 
such as phase noise. 
In Fig.~\ref{fig:const}\subref{const_blstm}, after BLSTM predistortion, the constellation cloud is significantly reduced, demonstrating 
effective compensation.


\subsection{BER performance}

\begin{figure}[!t]
    \centering
    \includegraphics[width=0.88\linewidth]{ber}
    \caption{BER performance with and without I/Q imbalance predistortion over a Rician fading channel}
    \label{ber}
\end{figure}

Fig. \ref{ber} shows the BER performance of two proposed predistortion schemes in a Rician channel with \textit{K} factors of 0 and 4.
For comparison, the results for the uncompensated case and the ideal case (no I/Q imbalance) are also plotted.
For $K=0$ at BER of $10^{-2}$, the RF signal-based and BLSTM-based schemes achieve approximately 6.0 dB and 5.5 dB performance improvements, 
respectively. 
For $K=4$ at BER of $10^{-5}$, the improvements are approximately 5.0 dB and 4.5 dB, respectively, compared to the uncompensated case.

As shown in Fig. \ref{ber}, the RF signal-based scheme shows near-ideal performance.
The BLSTM-based scheme exhibits about 0.5 dB performance gap compared to the RF signal-based scheme because the RF scheme directly measures errors in real time, while BLSTM relies on the training set covering limited ranges.
Moreover, the presence of residual MSE error $1.2 \times 10^{-4}$ and the computational complexity constraints that limit the model capacity result in the 0.5 dB performance gap.
However, the performance of the BLSTM-based predistortion scheme could be further improved by using more refined training set and the application of optimal layers and parameters.

These experimental results confirm that the proposed BLSTM architecture is effective for I/Q imbalance predistortion and achieves performance suitable for practical system implementation under the considered conditions.


\section{Conclusion}

This paper proposed two novel predistortion schemes for I/Q imbalance mitigation in OFDM systems.
These two predistortion schemes overcome the limitations of conventional techniques that require additional pilot signals and employ post-compensation with relatively large residual errors, providing an effective solution for mitigating I/Q imbalance exhibiting non-linear characteristics.

The performance of the proposed predistortion schemes is evaluated by simulation.
The simulation results show that the RF signal-based predistortion scheme achieves performance nearly identical to the ideal case without I/Q imbalance, whereas the BLSTM-based predistortion scheme shows a performance degradation of about 0.5 dB.
The remaining 0.5 dB gap mainly stems from limited training dataset coverage and model capacity constraints.
With further improvements in training strategies and model design, the performance gap between the proposed BLSTM-based predistortion scheme and the RF signal-based scheme may be further reduced.

%\section*{Acknowledgments}
%This should be a simple paragraph before the References to thank those %individuals and institutions who have supported your work on this article.

% \textit{et al.}

\begin{thebibliography}{1}
\bibliographystyle{IEEEtran}

\bibitem{REF1}
3GPP TS38.211, {\it{3rd Generation Partnership Project: 5G; NR; Physical channels and modulation}}. ETSI, Sophia Antipolis, France, 2018.

\bibitem{RP-25xxxx}
3GPP TSG RAN meeting \#109, {\it{Status Report to TSG}}. Beijing, China, Sep. 2025.

\bibitem{RP-252912}
3GPP TSG RAN meeting \#109, {\it{Study on 6G Radio}}. Beijing, China, Sep. 2025.

\bibitem{Prasad}
R. Prasad, {\it{OFDM for wireless communications systems}}, Norwood, MA, USA: Artech House, 2004.

\bibitem{REF4}
S. Woo, D. Lee, K. Kim, Y. Hur, C. Lee, and J. Laskar, ``Combined effects of RF impairments in the future IEEE 802.11n WLAN systems,'' in \textit{Proc. IEEE VTC-SPRING}, Stockholm, Sweden, 2005, pp. 1346-1349.

\bibitem{REF5}
H. Ryu, Y. S. Li, and J. Park, ``Nonlinear analysis of the phase noise in the OFDM communication system,'' \textit{IEEE Trans. Consumer Electron.}, vol. 50, no. 1, pp. 54-63, Feb. 2004.

\bibitem{REF6}
A. A. Abidi, ``Direct-conversion radio transceivers for digital communications,'' \textit{IEEE Journal of Solid-State Circuits}, vol. 30, No. 12, pp. 1399-1410, Dec. 1995.

\bibitem{Razavi}
B. Razavi, {\it{RF Microelectronics}}, Upper Saddle River, NJ, USA: Prentice Hall, 1998.

\bibitem{REF7}
A. Schuchert, R. Hasholzner, and P. Antonie, ``A novel IQ imbalance compensation scheme for the reception of OFDM signals,'' \textit{IEEE Trans. Consumer Electron.}, vol. 47, no. 3. pp. 313-318, Aug. 2001.

\bibitem{REF8}
L. Anttila, M. Valkama, and M. Renfors, ``Blind Compensation of Frequency-Selective I/Q Imbalances in Quadrature Radio Receivers: Circularity-Based Approach,'' in \textit{Proc. IEEE ICASSP}, Honolulu, USA, 2007, pp. 245-248.

\bibitem{REF9}
A. Kiayani, L. Anttila, Y. Zou, and M. Valkama, ``Advanced Receiver Design for Mitigating Multiple RF Impairments in OFDM Systems: Algorithms and RF Measurements,'' \textit{Journal of Electron. and Computer Engineer.}, vol. 2012, no. 3, pp. 1-16, Jan. 2012.

\bibitem{REF10}
Y. Zhu, C. Gong, J. Luo, M. Jin, X. Jin, and Z. Xu. ``Indoor Non-Line of Sight Visible Light Communication with a Bi-LSTM Neural Network,'' In \textit{Proc. 2020 IEEE Intern. Conf. on Commun. Workshops (ICC Workshops)}, Dublin, Ireland, 2020, pp. 1-6.

\bibitem{REF11}
Md. A. I. Sunny, M. M. S. Maswood, and A. G. Alharbi, ``Deep learning-based stock price prediction using LSTM and bi-directional LSTM model,'' In \textit{Proc. 2020 2nd Novel Intelligent and Leading Emerging Sciences Conference (NILES)}, Giza, Egypt, 2020, pp. 87-92.

\bibitem{REF12}
Y. Imrana, Y. Xiang, L. Ali, and Z. Abdul-Rauf. ``A bidirectional LSTM deep learning approach for intrusion detection,'' \textit{Expert Syst. with App.} vol. 185, no. 8, Dec. 2021.

\end{thebibliography}

\vfill

\end{document} 
